% =====================================================================
% chapters/07_Implications.tex — LaTeX rendering of ../sections/07_implications.en.md (2026-09-25 21:22:31)
% Section 7 of the SHORT paper, the last one. \input from main.tex, after 06_Non_tabulated.tex.
% This file DEFINES \label{sec:implications}, \label{sec:cost-of-deciding} (7.1),
%   \label{sec:limitations} (7.2) and \label{sec:conclusion} (7.3); no other chapter
%   references them.
% It REFERENCES sec:mechanism (Section 6.2, from 7.1) and sec:nontabulated (Section 6, from
%   7.2). sec:individual (Section 5.4) is no longer referenced from here.
% Figures: none, and none is wanted. ../PLAN.md § 7.1 forbids a cascade figure with empty
%   flow shares, and forbids \todo.
% Tables: none.
% Citation: chopra2024limits (AAMAS 2025), the one published answer to the inference cost. It
%   resolves in sample.bib.
% Section 7.2 carries four run-in limitations whose titles are bold in the master; they are
%   set with \textbf{} at the head of each paragraph.
% CUT CANDIDATE, from the master's comments, if the page budget requires it: the last sentence
%   of the cost paragraph of 7.3, "Billions of tokens per simulated day cannot be justified
%   just to reproduce what the survey already tabulates."
% ONE STATEMENT DEPENDS ON TICKET 103 although no placeholder is written here: the sentence
%   of the conclusion on the typed classifier reaching the same range (written under scenario
%   1; it leaves the conclusion under scenario 3).
% Re-carried by hand on 25 September 2026 against the master of 2026-09-25 21:22:31: 7.1
%   retitled ("What it costs to let generative agents decide", label renamed from
%   sec:architecture-implication, which nothing referenced); the three-stage division of
%   labour and the nominal-stage paragraph are withdrawn, replaced by two paragraphs (calling
%   generative agents only when an event occurs; which parts of the memory need a language
%   model, moved from 6.2); 7.2 rewritten as four titled limitations; 7.3 rewritten in five
%   paragraphs. Earlier states of this chapter live in git history.
% Maintained by hand (no generator since 2026-09-23): carry every validated pass of
%   ../sections/07_implications.en.md and set the date of line 2 to its marker.
% =====================================================================

\section{Implications, limitations, conclusion}
\label{sec:implications}

\subsection{What it costs to let generative agents decide}
\label{sec:cost-of-deciding}

With the models used, the cost of one simulated day determines where deliberation can be
afforded. One weekday for 1{,}000 personas requires 2{,}108 calls to the language model and
consumes 3 million tokens, with memory disabled. The 702 trips with a single option never
reach the model, and the measured run reused part of an earlier day. Turning memory on adds
some 2.5 million tokens more. Scaled to the whole study area (1.32 million residents aged five
and over), one simulated day would require 2.9 million calls, for 2.6 to 4.2 billion
tokens. One published answer to that cost asks the model once per behavioural archetype rather
than once per agent \citep{chopra2024limits}. There, 8.4 million agents cost some 400
queries. Agents of the same archetype share an estimated probability, and each draws its own
action from it. This saving does not apply here: archetypes are defined by attributes, and an
off-survey event is precisely what no attribute records.

At city scale, one way to afford generative agents would be to call them only when an event
occurs, and let a cheaper decision-maker handle the ordinary day. Our system does not include
this, and we do not estimate the share of trips that would reach the generative agent.

Such a split would also have to decide which parts of the memory need a language model. When
the model rates the severity of an event (Section~\ref{sec:mechanism}), it only picks one of
five levels, a label that a typed classifier could return as well. Recording the trace of what
happened and maintaining the belief drawn from it, by contrast, take written prose. Whether a
classifier could also do this is an open question.

\subsection{Limitations}
\label{sec:limitations}

Four limitations constrain what these measurements support.

\textbf{A local survey behind every step.} Prompt tuning, the training of the reference models
and the scoring of the benchmark all rest on a local travel survey. Toulouse has one,
available to researchers on request. Such surveys are collected very unevenly across
countries. Where none exists, prompts cannot be tuned, and their error on that territory
stays unknown. Generative agents therefore do not spare a territory the need for local
calibration data.

\textbf{A sampled and reported reference.} The reference is a sample of trips that residents
report the day after, from memory. It therefore contains a sampling error, largest in the
smallest strata, and misses the trips respondents forget. Benchmark scores measure a distance
to these reported shares, not to the trips actually made.

\textbf{One weekday as the reference.} The survey describes a single weekday outside school
holidays, for residents aged five and over, and our results hold for that day alone. Extending
the benchmark to weekends or holidays would need a reference for their activities, which differ
from those of a working day.

\textbf{No observed reaction to compare with.} Our survey records no reaction to an incident
or to an article. Section~\ref{sec:nontabulated} therefore shows that an agent reacts to an
event none of the 21 variables records. It does not show that the agent reacts as a resident
would.

\subsection{Conclusion}
\label{sec:conclusion}

We asked what verbalised deliberation adds to urban simulation, and where it earns its place
in a mobility agent. Published evaluations of generative mobility agents rarely confront them
with a real population. We confronted generative agents with a field survey, the 2023
household travel survey of the Toulouse area, in France. We placed them between simple
baselines and four reference models fitted on that survey, all reading the same 21 variables,
then pushed them as far as prompt tuning allowed.

On the ordinary day the survey describes, deliberation yields no measurable gain. A tuned
generative agent comes within the cohort resolution of the reference models, without
outperforming them. A typed classifier that writes no text reaches the same range for a
fiftieth of the cost. Without reference models on the same survey, none of these three
positions could have been read.

The modal split does not tell the whole story. The tuned agent and the gradient boosting model
score within the cohort resolution of each other, yet pick a different mode on one trip in
three. The typed classifier gets more reported trips wrong than a baseline that sends everyone
by car. An evaluation that scores the modal split alone, or one mode per trip, cannot see this
gap between shares and decisions.

The cost of deliberation grows with the population it simulates. One call per decision stays
affordable for a few dozen travellers, whereas a city multiplies the bill, in money and in
energy. A smaller model served on local hardware could bring that bill within reach, and later
language models may close the remaining gap. Whatever the model, what a decision-maker
consumes matters as much as its accuracy when choosing an architecture. Billions of tokens per
simulated day cannot be justified just to reproduce what the survey already tabulates.

Deliberation is worth its cost where the survey is silent. We traced one agent end to end over
several weeks. An engine failure entered its memory and cut its car use for fifteen days,
although the car stayed on offer, then faded.
